\documentclass[aspectratio=169]{beamer}
\usetheme{Madrid}
\usecolortheme{default}

\usepackage{amsmath,amssymb,amsfonts,bm}
\usepackage{physics}
\usepackage{mathtools}

\title{Summary of Five QAOA Gradient Methods}
\subtitle{Beamer Slides Adapted from the Uploaded Summary}
\author{}
\date{}

\begin{document}

\frame{\titlepage}

\begin{frame}{Outline}
\tableofcontents
\end{frame}

\section{Introduction}

\begin{frame}{Purpose of the document}
These slides summarize five methods for computing gradients in the
Quantum Approximate Optimization Algorithm (QAOA).

\medskip

Main goals:
\begin{itemize}
\item compare the formulas behind the different gradient estimators,
\item explain their implementation strategies,
\item analyze computational cost,
\item discuss practical tradeoffs between accuracy and hardware compatibility.
\end{itemize}

The motivating example is an XXZ Hamiltonian with
\[
L=6,\qquad p=2,\qquad N=2p=4
\]
variational parameters.
\end{frame}

\section{The QAOA Gradient Problem}

\begin{frame}{QAOA overview}
QAOA prepares a variational state of the form
\[
\ket{\psi(\bm{\gamma},\bm{\beta})}
=
U_M(\beta_p)U_C(\gamma_p)\cdots U_M(\beta_1)U_C(\gamma_1)\ket{+}^{\otimes n},
\]
where
\[
U_C(\gamma)=e^{-i\gamma H_C},
\qquad
U_M(\beta)=e^{-i\beta H_M}.
\]

Here:
\begin{itemize}
\item \(H_C\) is the cost Hamiltonian,
\item \(H_M\) is the mixer Hamiltonian,
\item \(p\) is the circuit depth,
\item \(\bm{\gamma},\bm{\beta}\) are variational parameters.
\end{itemize}
\end{frame}

\begin{frame}{Energy objective}
The variational objective is the energy expectation value
\[
E(\bm{\gamma},\bm{\beta})
=
\bra{\psi(\bm{\gamma},\bm{\beta})}H_C\ket{\psi(\bm{\gamma},\bm{\beta})}.
\]

If the cost Hamiltonian has a Pauli decomposition
\[
H_C = \sum_{\alpha=1}^{K} h_\alpha P_\alpha,
\]
then
\[
E(\bm{\gamma},\bm{\beta})
=
\sum_{\alpha=1}^{K} h_\alpha
\bra{\psi(\bm{\gamma},\bm{\beta})}P_\alpha\ket{\psi(\bm{\gamma},\bm{\beta})}.
\]
Thus the energy is a weighted sum of Pauli-string expectation values.
\end{frame}

\begin{frame}{Gradient challenge}
Optimization of the \(2p\) parameters requires the gradient
\[
\nabla E
=
\left(
\frac{\partial E}{\partial \gamma_1},
\frac{\partial E}{\partial \beta_1},
\dots,
\frac{\partial E}{\partial \gamma_p},
\frac{\partial E}{\partial \beta_p}
\right).
\]

For each parameter \(\theta_k\in\{\gamma_i,\beta_i\}\),
\[
\frac{\partial E}{\partial \theta_k}
=
\frac{\partial}{\partial \theta_k}
\bra{\psi(\bm{\theta})}H_C\ket{\psi(\bm{\theta})}
=
\sum_{\alpha=1}^{K} h_\alpha
\frac{\partial}{\partial \theta_k}\expval{P_\alpha}.
\]

The computational question is therefore:
\begin{quote}
How can one evaluate \(\partial_{\theta_k}\expval{P_\alpha}\) efficiently on quantum hardware?
\end{quote}
\end{frame}

\begin{frame}{General derivative structure}
Let \(\ket{\psi_0}=\ket{+}^{\otimes n}\) and let \(U(\bm{\theta})\) denote the full QAOA unitary.
Then
\[
E(\bm{\theta})
=
\bra{\psi_0}U^\dagger(\bm{\theta})H_CU(\bm{\theta})\ket{\psi_0}.
\]

Therefore
\[
\frac{\partial E}{\partial \theta_k}
=
\frac{\partial}{\partial \theta_k}
\bra{\psi_0}U^\dagger(\bm{\theta})H_CU(\bm{\theta})\ket{\psi_0}.
\]

If \(\theta_k\) parameterizes a gate
\[
U_k(\theta_k)=e^{-i\theta_k B_k},
\]
with generator \(B_k\in\{H_C,H_M\}\), and
\[
U = U_{\text{after}}\, U_k(\theta_k)\, U_{\text{before}},
\]
then all gradient methods reduce to evaluating derivatives of expectation values of the form
\[
\expval{U^\dagger B U}.
\]
\end{frame}

\begin{frame}{Why the problem is nontrivial}
The summary emphasizes several sources of difficulty:
\begin{enumerate}
\item Quantum hardware cannot directly apply \(H\ket{\psi}\) as a primitive operation.
\item Standard parameter-shift or finite-difference approaches require \(O(p)\) circuit evaluations per gradient.
\item Realistic Hamiltonians decompose into many Pauli terms
\[
H=\sum_\alpha h_\alpha P_\alpha,
\]
often with \(K\gg p\).
\item Non-uniform coefficients \(h_\alpha\) can degrade uniform-shift approximations.
\end{enumerate}
So the practical tradeoff is accuracy versus measurement/circuit cost.
\end{frame}

\begin{frame}{Test case: XXZ spin chain}
The summary uses the XXZ Hamiltonian
\[
H_{\mathrm{XXZ}}
=
\sum_{i=1}^{L-1}
\Big[
J_{xx}(\sigma_i^x\sigma_{i+1}^x+\sigma_i^y\sigma_{i+1}^y)
+J_z \sigma_i^z\sigma_{i+1}^z
\Big]
+
\sum_{i=1}^{L} h_z \sigma_i^z.
\]

For the example:
\[
J_{xx}=1.0,\qquad J_z=1.5,\qquad h_z=0.2,
\]
with
\[
L=6,\qquad p=2,\qquad N=4.
\]

This gives moderately non-uniform couplings with ratio
\[
h_{\max}/h_{\min}=1.5/0.2=7.5.
\]
\end{frame}

\section{Method 1: Exact Gradient}

\begin{frame}{Method 1: Exact gradient from Schr\"odinger-equation structure}
The exact gradient formula is given as
\[
\frac{\partial E}{\partial \theta_k}
=
i\bra{\psi}H_{\text{cost}}U^\dagger B_k U\ket{\psi}
-
i\bra{\psi}U^\dagger B_k U H_{\text{cost}}\ket{\psi}.
\]

Here:
\begin{itemize}
\item \(\ket{\psi}\) is the final QAOA state,
\item \(U\) is the full circuit unitary,
\item \(B_k\) is the generator corresponding to parameter \(\theta_k\).
\end{itemize}

This expression is exact and arises from differentiating the unitary evolution directly.
\end{frame}

\begin{frame}{Method 1: implementation and cost}
Implementation idea:
\begin{itemize}
\item apply the Hamiltonian MPO to intermediate states,
\item compute overlaps explicitly.
\end{itemize}

Computational cost:
\[
N \text{ circuit evaluations} + N \text{ Hamiltonian applications}.
\]

Assessment:
\begin{itemize}
\item \textbf{Pros:} exact gradient, minimal approximation error.
\item \textbf{Cons:} requires direct application of \(H\ket{\psi}\), which is generally not implementable on standard quantum hardware.
\end{itemize}
\end{frame}

\section{Method 2: Standard Parameter Shift}

\begin{frame}{Method 2: standard parameter-shift rule}
The standard parameter-shift philosophy is to evaluate
\[
\frac{\partial E}{\partial \theta_k}
\]
through shifted circuit parameters rather than direct differentiation.

For generators with suitable spectra, one obtains a formula of the general form
\[
\frac{\partial E}{\partial \theta_k}
=
c\Big(E(\theta_k+s)-E(\theta_k-s)\Big),
\]
for an appropriate shift \(s\) and prefactor \(c\).

This is the most common practical gradient method in variational quantum algorithms.
\end{frame}

\begin{frame}{Method 2: practical interpretation}
Method 2 is attractive because:
\begin{itemize}
\item it is hardware compatible,
\item it requires only expectation-value measurements of shifted circuits,
\item it avoids direct Hamiltonian-vector products.
\end{itemize}

However:
\begin{itemize}
\item it treats the generator more uniformly,
\item it may not exploit Hamiltonian structure,
\item it can become inaccurate or expensive when the Pauli decomposition is highly non-uniform.
\end{itemize}

The summary identifies this as the standard baseline method.
\end{frame}

\section{Method 3: Per-Pauli Parameter Shift}

\begin{frame}{Method 3: per-Pauli parameter shifts}
Instead of shifting the whole Hamiltonian generator uniformly, Method 3 decomposes
\[
B = \sum_\alpha h_\alpha P_\alpha
\]
and applies shifts at the level of the individual Pauli terms.

The idea is that different coefficients \(h_\alpha\) suggest different effective shifts, so one can adapt the parameter-shift rule to the structure of the Hamiltonian term by term.
\end{frame}

\begin{frame}{Method 3: advantages and drawbacks}
Potential advantages:
\begin{itemize}
\item improved accuracy for Hamiltonians with non-uniform coefficients,
\item finer control over gradient estimation,
\item more faithful treatment of heterogeneous Pauli decompositions.
\end{itemize}

Potential drawbacks:
\begin{itemize}
\item more circuit evaluations than uniform shifting,
\item higher implementation complexity,
\item more demanding bookkeeping at the Pauli-term level.
\end{itemize}
\end{frame}

\section{Method 4: Commuting-Group Optimization}

\begin{frame}{Method 4: commuting-group optimization}
Method 4 improves on Method 3 by grouping commuting Pauli terms together.

If a set of Pauli operators commute, then they can often be measured in a shared basis or shifted collectively.
This reduces redundant evaluations while retaining some of the accuracy advantages of per-term methods.
\end{frame}

\begin{frame}{Method 4: why commuting groups help}
Grouping commuting terms provides:
\begin{itemize}
\item reduced measurement overhead,
\item fewer effective shifted evaluations,
\item a more scalable route to per-Pauli methods.
\end{itemize}

The summary emphasizes that this can reduce cost significantly---often by about \(60\%\) to \(80\%\) compared with na\"ive per-Pauli strategies.
\end{frame}

\section{Method 5: Adaptive Hybrid Strategy}

\begin{frame}{Method 5: adaptive hybrid gradient strategy}
Method 5 combines ideas from the previous methods:
\begin{itemize}
\item use exact or refined shifts where Hamiltonian structure matters most,
\item use cheaper approximations where the effect of individual terms is small,
\item balance accuracy and circuit cost adaptively.
\end{itemize}

Thus Method 5 is intended as a practical middle ground between exactness and hardware efficiency.
\end{frame}

\begin{frame}{Method 5: motivation}
The motivation is clear:
\begin{itemize}
\item Method 1 is exact but impractical,
\item Method 2 is practical but may miss Hamiltonian-specific structure,
\item Methods 3 and 4 are more accurate but can be costly,
\item Method 5 seeks the best accuracy-cost compromise.
\end{itemize}

This makes it especially relevant when some Pauli terms dominate the gradient contribution more strongly than others.
\end{frame}

\section{Comparison and Tradeoffs}

\begin{frame}{Comparative questions posed by the summary}
The document frames the comparison through three main questions:
\begin{enumerate}
\item When does Hamiltonian structure matter?
\item What is the accuracy--cost tradeoff?
\item Can hybrid methods bridge the gap between theory and practice?
\end{enumerate}

These are exactly the right questions when deciding which QAOA gradient method to use in an implementation.
\end{frame}

\begin{frame}{Accuracy--cost picture}
The five methods fall into a rough hierarchy:

\begin{itemize}
\item \textbf{Method 1:} exact, low circuit count, but not hardware realizable.
\item \textbf{Method 2:} standard, simple, hardware compatible.
\item \textbf{Method 3:} more accurate for non-uniform Hamiltonians, but more expensive.
\item \textbf{Method 4:} recovers much of Method 3's benefit at reduced cost.
\item \textbf{Method 5:} adaptive compromise between accuracy and overhead.
\end{itemize}
\end{frame}

\begin{frame}{Concrete XXZ example}
For the XXZ chain with
\[
L=6,\qquad p=2,\qquad N=4,
\]
and approximately
\[
K \approx 14
\]
Pauli terms per Hamiltonian, the summary finds:

\begin{itemize}
\item non-uniform couplings make Methods 3--5 potentially more accurate than Method 2,
\item commuting-group and adaptive optimizations reduce cost substantially,
\item Method 2 remains the easiest standard implementation,
\item Methods 3--5 become attractive when the Hamiltonian coefficients vary strongly.
\end{itemize}
\end{frame}

\begin{frame}{Key insights from the summary}
The final conclusions highlight:

\begin{itemize}
\item For non-uniform coupling strengths, Methods 3--5 can outperform the standard parameter-shift approach in accuracy.
\item Methods 4 and 5 reduce the cost of per-Pauli approaches by roughly \(60\%\) to \(80\%\).
\item Method 2 remains the most common baseline because of simplicity and hardware compatibility.
\item The main novelty lies in synthesizing known shift and grouping ideas into QAOA-specific gradient strategies.
\end{itemize}
\end{frame}

\section{Novel Contributions}

\begin{frame}{Novel contributions claimed in the summary}
The document identifies four main contributions:
\begin{enumerate}
\item systematic application of per-Pauli parameter shifts with coefficient-dependent shifts to QAOA gradients,
\item a commuting-group optimization strategy,
\item an adaptive hybrid strategy balancing term importance and cost,
\item an empirical comparison showing when these methods improve on the standard approach.
\end{enumerate}

So the work is positioned as a novel synthesis rather than a completely new underlying gradient formalism.
\end{frame}

\section{Conclusions}

\begin{frame}{Overall conclusion}
The essential message is:
\begin{quote}
There is no single best QAOA gradient method for all Hamiltonians.
\end{quote}

Instead:
\begin{itemize}
\item hardware constraints favor simple parameter-shift methods,
\item structured Hamiltonians can justify more refined per-Pauli or grouped strategies,
\item adaptive hybrids may offer the best practical compromise.
\end{itemize}

Thus the ``best'' method depends on the geometry and decomposition of the Hamiltonian as well as the implementation platform.
\end{frame}

\begin{frame}{Possible extensions}
Natural follow-up topics include:
\begin{itemize}
\item explicit derivation of the standard parameter-shift rule for QAOA generators,
\item measurement grouping for larger Pauli decompositions,
\item shot-noise analysis for each method,
\item tensor-network or MPO simulation costs,
\item application to larger spin models and combinatorial optimization Hamiltonians.
\end{itemize}
\end{frame}

\begin{frame}{End}
\centering
{\Large Questions?}
\end{frame}

\end{document}
